\chapter{互斥集合}
\label{ch:disjoint}

% 本章的圖共用的樣式。 森林圖由模擬程式依虛擬碼逐步執行後產生，箭頭由子節點指向父節點。
\tikzset{
    dn/.style={circle, draw=black!65, minimum size=0.56cm, inner sep=0pt, font=\small},
    dnroot/.style={draw=coreStroke, very thick, fill=coreFill},   % 根（集合的代表）
    dnhl/.style={draw=memStroke, very thick, fill=memFill},       % 這一步父指標有改的節點
    dptr/.style={-{Stealth[length=1.6mm]}, draw=black!65, shorten >=0.5pt},
    plab/.style={font=\small},
    lnode/.style={draw=black!65, minimum width=0.62cm, minimum height=0.56cm, inner sep=1pt, font=\small},
    lset/.style={draw=coreStroke, thick, fill=coreFill, minimum height=0.56cm, inner sep=2pt, font=\small},
    lptr/.style={-{Stealth[length=1.6mm]}, draw=black!70},
    lrep/.style={-{Stealth[length=1.4mm]}, draw=coreStroke!80, densely dashed},
}

系上的電腦教室有幾十台電腦，網管這學期要重新拉網路線。 他手上有一張表，一行寫一
條線接了哪兩台電腦，照著表一條一條接。 接線的同時，助教隨時會跑來問：「第 7 台和第
23 台現在連得到嗎？」 線接上以後不會再拆，大家只在乎誰和誰在同一個網段。 如果每
接一條線，就把整間教室走一遍，替每台電腦重新標上網段編號，接幾百條線就要走幾百
遍。 有沒有辦法讓每次接線、每次回答都只做很少的工作？ 往下讀之前，先想想你會怎麼
記錄「誰和誰同一組」。

這一章的資料結構叫\textbf{互斥集合}（disjoint sets），也常叫做 union--find。 它只
支援三件事：建立只有一個元素的集合、問一個元素在哪個集合、把兩個集合合併。 功能很
少，卻是許多圖論演算法的零件，第~24~章的 Kruskal 最小生成樹演算法就靠它判斷一條邊
會不會形成迴路。 這一章先用串列表示集合，看「把短的接到長的後面」這個簡單的規則能省
下多少工作；再改用第~7~章說過的「只存父節點」的森林，配上依大小或依秩合併、路徑壓
縮兩個技巧，最後說明這兩個技巧合起來，每個操作的平攤成本幾乎是常數。

這一章用到《C 程式設計基礎》第~14~章的結構、第~15~章的遞迴和第~19~章的
\texttt{malloc}、\texttt{free}；平攤的意思見第~3~章，森林和父節點陣列見第~7~章。

本章的\textbf{高度}照第~7~章的定義數節點：只有一個節點的樹高度是 $1$，空樹是 $0$；
\textbf{深度}是從根往下走的邊數，根的深度是 $0$。 第~\ref{sec:ds-rank}~節會出現另
一個量叫\textbf{秩}，它是另外定義的，和高度差一，不要混在一起。

\section{操作與串列表示}
\label{sec:ds-list}

\subsection{三個操作}

互斥集合維護一群\textbf{兩兩沒有共同元素}的集合 $S_1, S_2, \ldots, S_k$。 每個集合選
一個成員當\textbf{代表}（representative），用代表來指稱整個集合。 代表是誰可以由實作
決定，但只要集合沒有被改過，同一個集合問兩次，必須得到同一個代表。 支援的操作有三
個：

\begin{itemize}
\item \proc{Make-Set}$(x)$：建立只含 $x$ 的新集合，$x$ 就是它的代表。 $x$ 原本不能在
任何集合裡。
\item \proc{Find-Set}$(x)$：回傳 $x$ 所在集合的代表。
\item \proc{Union}$(x, y)$：把 $x$ 所在的集合和 $y$ 所在的集合合成一個，原來的兩個集
合不再存在。 新集合的代表從兩個舊代表裡選一個。 $x$、$y$ 已經在同一個集合時什麼都
不做。
\end{itemize}

問「$x$ 和 $y$ 在不在同一組」，就是比較 $\proc{Find-Set}(x) = \proc{Find-Set}(y)$。
互斥集合不支援把一個集合拆開，也不支援把元素移出集合；第~\ref{sec:ds-app}~節會看到
遇到「拆線」時怎麼繞過這個限制。

分析時用兩個參數：$n$ 是 \proc{Make-Set} 的次數，也就是元素的總數；$m$ 是所有操作
的總數，三種都算。 每個元素都要先做一次 \proc{Make-Set}，所以 $m \ge n$。 每做一次
真正的合併，集合就少一個，所以真正的合併最多 $n - 1$ 次。

章首的教室用這三個操作來寫：每台電腦做一次 \proc{Make-Set}，每接一條線 $(u, v)$ 做一次
\proc{Union}$(u, v)$；助教問 $u$、$v$ 連不連得到，就回答 \proc{Find-Set}$(u) =
\proc{Find-Set}(v)$ 是否成立。 如果所有的線一次給齊，也可以用第~21~章的廣度或深度優
先搜尋，一次把每台電腦標上組別；互斥集合的好處在於線是一條一條來的，每來一條線只做一
次合併，中間隨時可以問。

\subsection{串列表示}

最直接的做法是把每個集合存成一條鏈結串列（第~5~章）。 每個集合有一個\textbf{集合物
件}，記錄串列的頭 $\mathit{head}$、尾 $\mathit{tail}$ 和元素個數 $\mathit{size}$；串
列裡每個元素 $x$ 除了往下一個的指標 $x.\mathit{next}$，還多一個指回集合物件的指標
$x.\mathit{set}$。 代表就是串列的第一個元素（圖~\ref{fig:ds-list}）。

\begin{center}
\begin{tikzpicture}[x=1cm, y=1cm]
\node[lset] (sa) at (1.4, 1.25) {$\mathit{head}\ \ \mathit{tail}\ \ \mathit{size}=3$};
\node[lnode] (a4) at (0, 0) {4};
\node[lnode] (a9) at (1.4, 0) {9};
\node[lnode] (a2) at (2.8, 0) {2};
\node[font=\small] (an) at (3.9, 0) {\NIL};
\draw[lptr] (a4) -- (a9);
\draw[lptr] (a9) -- (a2);
\draw[lptr] (a2) -- (an);
\foreach \v in {a4, a9, a2} { \draw[lrep] (\v.north) -- (\v.north |- sa.south); }
\node[plab] at (1.4, -0.85) {集合 $\{4, 9, 2\}$，代表是 $4$};
\begin{scope}[xshift=6.2cm]
\node[lset] (sb) at (0.7, 1.25) {$\mathit{head}\ \ \mathit{tail}\ \ \mathit{size}=2$};
\node[lnode] (b6) at (0, 0) {6};
\node[lnode] (b1) at (1.4, 0) {1};
\node[font=\small] (bn) at (2.5, 0) {\NIL};
\draw[lptr] (b6) -- (b1);
\draw[lptr] (b1) -- (bn);
\foreach \v in {b6, b1} { \draw[lrep] (\v.north) -- (\v.north |- sb.south); }
\node[plab] at (0.7, -0.85) {集合 $\{6, 1\}$，代表是 $6$};
\end{scope}
\end{tikzpicture}
\diagcaption{用串列表示兩個集合。 實線是 $\mathit{next}$，虛線是每個元素指回集合物件
的 $\mathit{set}$。 \proc{Find-Set}$(2)$ 沿虛線找到集合物件，再讀 $\mathit{head}$，
得到 $4$。}
\label{fig:ds-list}
\end{center}

\begin{pseudo}[label=alg:ds-list]{串列表示，加權合併}
\Procedure{Make-Set}{$x$}
    \State 建立集合物件 $S$；$S.\mathit{head} \gets x$，$S.\mathit{tail} \gets x$，$S.\mathit{size} \gets 1$
    \State $x.\mathit{set} \gets S$；$x.\mathit{next} \gets \NIL$
\EndProcedure
\Function{Find-Set}{$x$}
    \State \Return $x.\mathit{set}.\mathit{head}$
\EndFunction
\Procedure{Union}{$x, y$}
    \State $S \gets x.\mathit{set}$
    \State $T \gets y.\mathit{set}$
    \If{$S = T$}
        \State \Return
    \EndIf
    \If{$S.\mathit{size} < T.\mathit{size}$} \Comment{加權：讓 $T$ 是比較短的那條}
        \State 交換 $S$ 和 $T$
    \EndIf
    \State $S.\mathit{tail}.\mathit{next} \gets T.\mathit{head}$
    \State $S.\mathit{tail} \gets T.\mathit{tail}$
    \State $z \gets T.\mathit{head}$
    \While{$z \ne \NIL$} \Comment{$T$ 的每個元素改指 $S$}
        \State $z.\mathit{set} \gets S$
        \State $z \gets z.\mathit{next}$
    \EndWhile
    \State $S.\mathit{size} \gets S.\mathit{size} + T.\mathit{size}$
    \State 釋放集合物件 $T$
\EndProcedure
\end{pseudo}

\proc{Make-Set} 和 \proc{Find-Set} 都是 $O(1)$。 \proc{Union} 把 $T$ 整條接到 $S$ 的
尾巴，接上去只要改兩個指標（第 13--14 行），貴的是第 16--18 行：$T$ 的每個元素都要把
$\mathit{set}$ 改成 $S$，花 $\Theta(T.\mathit{size})$ 的時間。

先看拿掉第 11--12 行、永遠把 $y$ 的串列接到 $x$ 後面的版本。 教室裡新到一台電腦就
接到已經連好的大群組上，寫成操作是：\proc{Make-Set} $x_1, \ldots, x_n$，再依序做
\[
\proc{Union}(x_2, x_1),\ \proc{Union}(x_3, x_1),\ \ldots,\ \proc{Union}(x_n, x_1).
\]
第 $i$ 次合併時，$x_1$ 所在的集合已經有 $i$ 個元素，整條被接到只有一個元素的新集合
後面，要改 $i$ 個 $\mathit{set}$ 指標。 總共改了 $1 + 2 + \cdots + (n - 1) =
n(n-1)/2$ 次，而操作數只有 $m = 2n - 1$。 平均每個操作要 $\Theta(n)$ 的時間。

毛病出在每次都把長的接到短的後面。 第 11--12 行的\textbf{加權合併}（weighted-union
heuristic）規則是：比較兩條串列的長度，把短的接到長的後面，只改短的那一條的指標。
一樣長時接哪一條都可以，本章約定接 $y$ 的那一條。 同樣那串操作，每次只要改 $1$ 個
指標。

\begin{theorem}
\label{thm:ds-weighted-list}
串列表示使用加權合併時，$m$ 個操作（其中 $n$ 個是 \proc{Make-Set}）的總時間是
$O(m + n \lg n)$。
\end{theorem}

\begin{proof}
每個操作除了改 $\mathit{set}$ 指標以外，都只花常數時間，合計 $O(m)$。 剩下的是改指
標的總次數。 固定一個元素 $x$，看它的 $\mathit{set}$ 指標被改的時候：$x$ 一定在比較
短（或一樣長）的那一條裡，合併後的集合至少是 $x$ 原本所在集合的兩倍大。 $x$ 的集合
一開始只有 $1$ 個元素，被改過 $k$ 次之後至少有 $2^k$ 個元素；集合不會超過 $n$ 個元
素，所以 $2^k \le n$，$k \le \lg n$。 $n$ 個元素加起來最多改 $n \lg n$ 次，總時間
$O(m + n \lg n)$。
\end{proof}

這個上界在最壞情況是緊的。 設 $n = 2^k$，先兩兩合併成 $n/2$ 個大小 $2$ 的集合，再兩
兩合併成 $n/4$ 個大小 $4$ 的集合……一共 $k$ 輪。 每一輪的每次合併都是兩條一樣長的
串列，有一半的元素要改指標，每一輪改 $n/2$ 次，總共 $(n/2)\lg n$ 次。

\section{森林表示}
\label{sec:ds-forest}

串列表示讓 \proc{Find-Set} 很快，代價是合併要改一整條的指標。 反過來想：合併時只改
一個指標，\proc{Find-Set} 多走幾步。 做法是把每個集合存成一棵有根樹，每個節點只記
它的父節點 $p[x]$，根就是代表，根的父節點設成自己。 第~7~章說過，一個陣列
$p[1..n]$ 就能存下整個森林。

\begin{pseudo}[label=alg:ds-naive]{森林表示，不用任何技巧}
\Procedure{Make-Set}{$x$}
    \State $p[x] \gets x$
\EndProcedure
\Function{Find-Set}{$x$}
    \While{$p[x] \ne x$}
        \State $x \gets p[x]$
    \EndWhile
    \State \Return $x$
\EndFunction
\Procedure{Union}{$x, y$}
    \State $r_x \gets \Call{Find-Set}{x}$
    \State $r_y \gets \Call{Find-Set}{y}$
    \If{$r_x \ne r_y$}
        \State $p[r_y] \gets r_x$ \Comment{$y$ 的根接到 $x$ 的根底下}
    \EndIf
\EndProcedure
\end{pseudo}

第 11 行接的是兩個\textbf{根}。 寫成 $p[y] \gets x$ 是常見的錯：$y$ 原本的根和它其他
的子孫會留在原來的樹裡，兩個集合沒有真的合起來。

\proc{Find-Set} 從 $x$ 走到根，時間和 $x$ 的深度成正比；\proc{Union} 除了兩次
\proc{Find-Set} 只有常數時間。 問題是樹可能長得很高。 把第~\ref{sec:ds-list}~節那串
操作搬過來：\proc{Union}$(x_{i+1}, x_1)$ 每次都把 $x_1$ 所在的樹的根接到新節點底下，
樹長成一條 $x_1 \to x_2 \to \cdots \to x_n$ 的鏈，第 $i$ 次合併光是
\proc{Find-Set}$(x_1)$ 就要經過 $i$ 個節點，總時間又是 $\Theta(n^2)$。 下面兩個技巧
分別從「合併時怎麼接」和「找的時候順便整理」下手。

\subsection{依大小合併}
\label{sec:ds-size}

合併時讓節點比較少的那棵樹的根，接到節點比較多的那棵樹的根底下，叫做\textbf{依大小
合併}（union by size）。 每個根多記一個 $\mathit{size}$，是這棵樹的節點數；不是根的
節點的 $\mathit{size}$ 不再使用。

\begin{pseudo}[label=alg:ds-size]{依大小合併（\proc{Make-Set}$(x)$ 另外設 $\mathit{size}[x] \gets 1$）}
\Procedure{Union}{$x, y$}
    \State $r_x \gets \Call{Find-Set}{x}$
    \State $r_y \gets \Call{Find-Set}{y}$
    \If{$r_x = r_y$}
        \State \Return
    \EndIf
    \If{$\mathit{size}[r_x] < \mathit{size}[r_y]$}
        \State 交換 $r_x$ 和 $r_y$ \Comment{讓 $r_x$ 是節點比較多的根}
    \EndIf
    \State $p[r_y] \gets r_x$
    \State $\mathit{size}[r_x] \gets \mathit{size}[r_x] + \mathit{size}[r_y]$
\EndProcedure
\end{pseudo}

第 4--5 行不能省：$x$、$y$ 已經同一組時如果照樣執行第 9 行，$\mathit{size}$ 會被重
複累加。 兩棵樹一樣大時第 6 行不成立，$y$ 的根接到 $x$ 的根底下，和
虛擬碼~\ref{alg:ds-naive} 的方向相同。 第~\ref{sec:ds-cost}~節會證明，依大小合併的
樹，高度不超過 $\lfloor \lg n \rfloor + 1$。

\subsection{依秩合併}
\label{sec:ds-rank}

另一種接法是比較兩棵樹的\textbf{秩}（rank）。 秩是每個節點記的一個非負整數，規則只
有三條：

\begin{itemize}
\item \proc{Make-Set}$(x)$ 時 $\mathit{rank}[x] = 0$。 只有一個節點的樹，根的秩是
$0$（這棵樹的高度是 $1$）。
\item 合併兩個秩不同的根，秩小的接到秩大的底下，兩個秩都不變。
\item 合併兩個秩相同的根，$y$ 的根接到 $x$ 的根底下，留下來的根秩加一。
\end{itemize}

\begin{pseudo}[label=alg:ds-rank]{依秩合併（\proc{Make-Set}$(x)$ 另外設 $\mathit{rank}[x] \gets 0$）}
\Procedure{Union}{$x, y$}
    \State $r_x \gets \Call{Find-Set}{x}$
    \State $r_y \gets \Call{Find-Set}{y}$
    \If{$r_x = r_y$}
        \State \Return
    \EndIf
    \If{$\mathit{rank}[r_x] < \mathit{rank}[r_y]$}
        \State 交換 $r_x$ 和 $r_y$
    \EndIf
    \State $p[r_y] \gets r_x$
    \If{$\mathit{rank}[r_x] = \mathit{rank}[r_y]$}
        \State $\mathit{rank}[r_x] \gets \mathit{rank}[r_x] + 1$
    \EndIf
\EndProcedure
\end{pseudo}

秩和高度的關係要分兩種情況說。 不做路徑壓縮時，節點 $x$ 的秩\textbf{恰好}等於以 $x$
為根的子樹裡，從 $x$ 往下最長的路徑有幾條邊（下一節的路徑壓縮會讓樹變矮，
但不改任何秩）。 所以一般來說，秩是這個邊數的上界。 換成本書數節點的高度：
\[
\text{以 $x$ 為根的子樹的高度} \le \mathit{rank}[x] + 1,
\]
不做路徑壓縮時等號成立。 秩只在根上會變，一個節點變成別人的子節點以後，它的秩就固
定了。

\begin{examnote}
CLRS 和許多考卷用「邊數」定義高度，單一節點的高度是 $0$，所以會看到「秩是高度的上
界」「不壓縮時秩等於高度」這種寫法；換成本書數節點的高度，要寫成高度 $\le$ 秩 $+ 1$。
也有書把依秩合併叫做\textbf{依高度合併}（union by height），內容相同。 作答時先寫出
自己用的定義，題目問「秩」就照秩的三條規則算，不要拿實際的樹高代替。
\end{examnote}

依大小和依秩多數時候做出一樣的選擇，但不一定。 例如一棵是 $5$ 個節點、高度 $2$ 的
星形（根的秩 $1$），另一棵是 $4$ 個節點、高度 $3$ 的樹（根的秩 $2$）：依大小合併會
讓星形的根留下來，依秩合併會讓另一棵的根留下來。 觀念題第~\ref{q:ds-compare}~題就
是這種情形。

\subsection{路徑壓縮}
\label{sec:ds-compress}

\proc{Find-Set}$(x)$ 已經走過從 $x$ 到根的整條路徑，順便把路上每個節點的父指標都直接
改成根，下次再找這些節點，一步就到。 這叫\textbf{路徑壓縮}（path compression）。 它
只改父指標，不改任何秩和 $\mathit{size}$。

\begin{pseudo}[label=alg:ds-find-rec]{路徑壓縮（遞迴）}
\Function{Find-Set}{$x$}
    \If{$p[x] \ne x$}
        \State $p[x] \gets \Call{Find-Set}{p[x]}$
    \EndIf
    \State \Return $p[x]$
\EndFunction
\end{pseudo}

遞迴版一路往上找到根，回來的時候（第 3 行）把根的編號一層一層填回去。 它其實走了兩
趟：去的時候找根，回來的時候改指標。 寫成迴圈就把兩趟分開：

\begin{pseudo}[label=alg:ds-find-iter]{路徑壓縮（迴圈）}
\Function{Find-Set-Iter}{$x$}
    \State $r \gets x$
    \While{$p[r] \ne r$} \Comment{第一趟：找根}
        \State $r \gets p[r]$
    \EndWhile
    \While{$p[x] \ne r$} \Comment{第二趟：路上的節點都改接到根}
        \State $\mathit{next} \gets p[x]$
        \State $p[x] \gets r$
        \State $x \gets \mathit{next}$
    \EndWhile
    \State \Return $r$
\EndFunction
\end{pseudo}

第 6 行要先記下 $p[x]$，第 7 行改掉之後就找不到原本的下一個節點了。 兩個版本對同一
個森林做同一次 \proc{Find-Set}，得到的父節點陣列完全一樣，第~\ref{sec:ds-c}~節的 C
程式兩種都用到。 遞迴版寫起來短，但遞迴的深度等於路徑的長度；如果不搭配依大小或依
秩合併，樹可以是一條長度 $n$ 的鏈，$n$ 很大時會用光呼叫堆疊。

\subsection{路徑減半與路徑分割}
\label{sec:ds-onepass}

路徑壓縮要走兩趟。 另外有兩種只走一趟的變形，不必先找到根：
\begin{itemize}
\item \textbf{路徑減半}（path halving）：路上每隔一個節點，把它的父指標改成祖父，再
跳到（新的）父節點。
\item \textbf{路徑分割}（path splitting）：路上\textbf{每個}節點都把父指標改成祖父，再
跳到原本的父節點。
\end{itemize}

\begin{pseudo}[label=alg:ds-find-onepass]{只走一趟的兩種變形}
\Function{Find-Halving}{$x$}
    \While{$p[x] \ne x$}
        \State $p[x] \gets p[p[x]]$
        \State $x \gets p[x]$
    \EndWhile
    \State \Return $x$
\EndFunction
\Function{Find-Splitting}{$x$}
    \While{$p[x] \ne x$}
        \State $\mathit{next} \gets p[x]$
        \State $p[x] \gets p[\mathit{next}]$
        \State $x \gets \mathit{next}$
    \EndWhile
    \State \Return $x$
\EndFunction
\end{pseudo}

圖~\ref{fig:ds-onepass} 是同一條路徑 $1 \to 2 \to \cdots \to 6$ 做
\proc{Find-Set}$(1)$ 的三種結果，橘色是父指標有改的節點。 路徑壓縮讓整條路徑都掛到根
底下；減半只改了 $1$ 和 $3$，路徑長度大約減半；分割改了 $1$ 到 $4$，原來的路徑拆成
$1 \to 3 \to 5 \to 6$ 和 $2 \to 4 \to 6$ 兩條。 一次 \proc{Find-Set} 之後，減半和分
割的樹都比壓縮的高，但 Tarjan 和 van Leeuwen 證明，搭配依秩（或依大小）合併時，三種
做法的總時間是同一級（第~\ref{sec:ds-cost}~節）。

\begin{center}
\begin{tikzpicture}[x=0.8cm,y=0.85cm,baseline=(current bounding box.north)]
\node[dn] (pa1) at (0,-5) {1};
\node[dn] (pa2) at (0,-4) {2};
\node[dn] (pa3) at (0,-3) {3};
\node[dn] (pa4) at (0,-2) {4};
\node[dn] (pa5) at (0,-1) {5};
\node[dn,dnroot] (pa6) at (0,0) {6};
\draw[dptr] (pa1) -- (pa2);
\draw[dptr] (pa2) -- (pa3);
\draw[dptr] (pa3) -- (pa4);
\draw[dptr] (pa4) -- (pa5);
\draw[dptr] (pa5) -- (pa6);
\node[plab, anchor=north] at (0,-5.45) {(a) 原來的路徑};
\end{tikzpicture}\hspace{0.8cm}%
\begin{tikzpicture}[x=0.8cm,y=0.85cm,baseline=(current bounding box.north)]
\node[dn,dnhl] (pb1) at (0,-1) {1};
\node[dn,dnhl] (pb2) at (1,-1) {2};
\node[dn,dnhl] (pb3) at (2,-1) {3};
\node[dn,dnhl] (pb4) at (3,-1) {4};
\node[dn] (pb5) at (4,-1) {5};
\node[dn,dnroot] (pb6) at (2,0) {6};
\draw[dptr] (pb1) -- (pb6);
\draw[dptr] (pb2) -- (pb6);
\draw[dptr] (pb3) -- (pb6);
\draw[dptr] (pb4) -- (pb6);
\draw[dptr] (pb5) -- (pb6);
\node[plab, anchor=north] at (2,-5.45) {(b) 路徑壓縮};
\end{tikzpicture}\hspace{0.8cm}%
\begin{tikzpicture}[x=0.8cm,y=0.85cm,baseline=(current bounding box.north)]
\node[dn,dnhl] (pc1) at (0,-3) {1};
\node[dn] (pc2) at (1,-3) {2};
\node[dn,dnhl] (pc3) at (0.5,-2) {3};
\node[dn] (pc4) at (2,-2) {4};
\node[dn] (pc5) at (1.25,-1) {5};
\node[dn,dnroot] (pc6) at (1.25,0) {6};
\draw[dptr] (pc1) -- (pc3);
\draw[dptr] (pc2) -- (pc3);
\draw[dptr] (pc3) -- (pc5);
\draw[dptr] (pc4) -- (pc5);
\draw[dptr] (pc5) -- (pc6);
\node[plab, anchor=north] at (1.25,-5.45) {(c) 路徑減半};
\end{tikzpicture}\hspace{0.8cm}%
\begin{tikzpicture}[x=0.8cm,y=0.85cm,baseline=(current bounding box.north)]
\node[dn,dnhl] (pd1) at (1,-3) {1};
\node[dn,dnhl] (pd2) at (0,-2) {2};
\node[dn,dnhl] (pd3) at (1,-2) {3};
\node[dn,dnhl] (pd4) at (0,-1) {4};
\node[dn] (pd5) at (1,-1) {5};
\node[dn,dnroot] (pd6) at (0.5,0) {6};
\draw[dptr] (pd1) -- (pd3);
\draw[dptr] (pd2) -- (pd4);
\draw[dptr] (pd3) -- (pd5);
\draw[dptr] (pd4) -- (pd6);
\draw[dptr] (pd5) -- (pd6);
\node[plab, anchor=north] at (0.5,-5.45) {(d) 路徑分割};
\end{tikzpicture}
\diagcaption{對同一條路徑做 \proc{Find-Set}$(1)$。 藍框是根，橘色是父指標有改的節點，
箭頭由子節點指向父節點。 父節點陣列 $p[1..6]$ 依序是 (b) $6,6,6,6,6,6$；(c)
$3,3,5,5,6,6$；(d) $3,4,5,6,6,6$。}
\label{fig:ds-onepass}
\end{center}

\subsection{手算追蹤：實驗室網路分組}
\label{sec:ds-trace}

實驗室有 $10$ 台電腦，編號 $1$ 到 $10$，一開始各自一組。 網管依序做了下面 $11$ 個操
作，用\textbf{依大小合併}（虛擬碼~\ref{alg:ds-size}）加\textbf{路徑壓縮}（虛擬
碼~\ref{alg:ds-find-rec}），一樣大時 $y$ 的根接到 $x$ 的根底下。 \proc{Union} 裡的兩
次 \proc{Find-Set} 也會壓縮路徑，這一點常被忽略。 下表每一列是做完那個操作之後的
$p[1..10]$，粗體是這一步改過的格子；圖~\ref{fig:ds-trace} 畫出其中四個時間點的森林。

\begin{center}
\small
\setlength{\tabcolsep}{4.2pt}
\begin{tabular}{rl*{10}{c}}
\toprule
步 & 操作 & $p[1]$ & $p[2]$ & $p[3]$ & $p[4]$ & $p[5]$ & $p[6]$ & $p[7]$ & $p[8]$ & $p[9]$ & $p[10]$ \\
\midrule
0 & 開始 & 1 & 2 & 3 & 4 & 5 & 6 & 7 & 8 & 9 & 10 \\
1 & \proc{Union}$(1, 2)$ & 1 & \textbf{1} & 3 & 4 & 5 & 6 & 7 & 8 & 9 & 10 \\
2 & \proc{Union}$(3, 4)$ & 1 & 1 & 3 & \textbf{3} & 5 & 6 & 7 & 8 & 9 & 10 \\
3 & \proc{Union}$(5, 6)$ & 1 & 1 & 3 & 3 & 5 & \textbf{5} & 7 & 8 & 9 & 10 \\
4 & \proc{Union}$(7, 8)$ & 1 & 1 & 3 & 3 & 5 & 5 & 7 & \textbf{7} & 9 & 10 \\
5 & \proc{Union}$(2, 4)$ & 1 & 1 & \textbf{1} & 3 & 5 & 5 & 7 & 7 & 9 & 10 \\
6 & \proc{Union}$(6, 8)$ & 1 & 1 & 1 & 3 & 5 & 5 & \textbf{5} & 7 & 9 & 10 \\
7 & \proc{Union}$(9, 10)$ & 1 & 1 & 1 & 3 & 5 & 5 & 5 & 7 & 9 & \textbf{9} \\
8 & \proc{Union}$(6, 2)$ & \textbf{5} & 1 & 1 & 3 & 5 & 5 & 5 & 7 & 9 & 9 \\
9 & \proc{Find-Set}$(4)$ & 5 & 1 & \textbf{5} & \textbf{5} & 5 & 5 & 5 & 7 & 9 & 9 \\
10 & \proc{Union}$(10, 8)$ & 5 & 1 & 5 & 5 & 5 & 5 & 5 & \textbf{5} & \textbf{5} & 9 \\
11 & \proc{Find-Set}$(10)$ & 5 & 1 & 5 & 5 & 5 & 5 & 5 & 5 & 5 & \textbf{5} \\
\bottomrule
\end{tabular}
\end{center}

\begin{center}
\begin{tikzpicture}[x=0.8cm,y=0.85cm,baseline=(current bounding box.north)]
\node[dn,dnroot] (tr6x1) at (0.5,0) {1};
\node[dn] (tr6x2) at (0,-1) {2};
\node[dn] (tr6x3) at (1,-1) {3};
\node[dn] (tr6x4) at (1,-2) {4};
\node[dn,dnroot] (tr6x5) at (2.85,0) {5};
\node[dn] (tr6x6) at (2.35,-1) {6};
\node[dn,dnhl] (tr6x7) at (3.35,-1) {7};
\node[dn] (tr6x8) at (3.35,-2) {8};
\node[dn,dnroot] (tr6x9) at (4.7,0) {9};
\node[dn,dnroot] (tr6x10) at (6.05,0) {10};
\draw[dptr] (tr6x2) -- (tr6x1);
\draw[dptr] (tr6x3) -- (tr6x1);
\draw[dptr] (tr6x4) -- (tr6x3);
\draw[dptr] (tr6x6) -- (tr6x5);
\draw[dptr] (tr6x7) -- (tr6x5);
\draw[dptr] (tr6x8) -- (tr6x7);
\node[plab, anchor=north] at (3.02,-3.45) {(a) 第 6 步之後};
\end{tikzpicture}\hspace{1.0cm}%
\begin{tikzpicture}[x=0.8cm,y=0.85cm,baseline=(current bounding box.north)]
\node[dn,dnhl] (tr8x1) at (0.5,-1) {1};
\node[dn] (tr8x2) at (0,-2) {2};
\node[dn] (tr8x3) at (1,-2) {3};
\node[dn] (tr8x4) at (1,-3) {4};
\node[dn,dnroot] (tr8x5) at (1.75,0) {5};
\node[dn] (tr8x6) at (2,-1) {6};
\node[dn] (tr8x7) at (3,-1) {7};
\node[dn] (tr8x8) at (3,-2) {8};
\node[dn,dnroot] (tr8x9) at (4.35,0) {9};
\node[dn] (tr8x10) at (4.35,-1) {10};
\draw[dptr] (tr8x1) -- (tr8x5);
\draw[dptr] (tr8x2) -- (tr8x1);
\draw[dptr] (tr8x3) -- (tr8x1);
\draw[dptr] (tr8x4) -- (tr8x3);
\draw[dptr] (tr8x6) -- (tr8x5);
\draw[dptr] (tr8x7) -- (tr8x5);
\draw[dptr] (tr8x8) -- (tr8x7);
\draw[dptr] (tr8x10) -- (tr8x9);
\node[plab, anchor=north] at (2.17,-3.45) {(b) 第 8 步之後};
\end{tikzpicture}

\vspace{0.5em}
\begin{tikzpicture}[x=0.8cm,y=0.85cm,baseline=(current bounding box.north)]
\node[dn] (tr9x1) at (0,-1) {1};
\node[dn] (tr9x2) at (0,-2) {2};
\node[dn,dnhl] (tr9x3) at (1,-1) {3};
\node[dn,dnhl] (tr9x4) at (2,-1) {4};
\node[dn,dnroot] (tr9x5) at (2,0) {5};
\node[dn] (tr9x6) at (3,-1) {6};
\node[dn] (tr9x7) at (4,-1) {7};
\node[dn] (tr9x8) at (4,-2) {8};
\node[dn,dnroot] (tr9x9) at (5.35,0) {9};
\node[dn] (tr9x10) at (5.35,-1) {10};
\draw[dptr] (tr9x1) -- (tr9x5);
\draw[dptr] (tr9x2) -- (tr9x1);
\draw[dptr] (tr9x3) -- (tr9x5);
\draw[dptr] (tr9x4) -- (tr9x5);
\draw[dptr] (tr9x6) -- (tr9x5);
\draw[dptr] (tr9x7) -- (tr9x5);
\draw[dptr] (tr9x8) -- (tr9x7);
\draw[dptr] (tr9x10) -- (tr9x9);
\node[plab, anchor=north] at (2.67,-2.45) {(c) 第 9 步之後};
\end{tikzpicture}\hspace{0.8cm}%
\begin{tikzpicture}[x=0.8cm,y=0.85cm,baseline=(current bounding box.north)]
\node[dn] (tr11x1) at (0,-1) {1};
\node[dn] (tr11x2) at (0,-2) {2};
\node[dn] (tr11x3) at (1,-1) {3};
\node[dn] (tr11x4) at (2,-1) {4};
\node[dn,dnroot] (tr11x5) at (3.5,0) {5};
\node[dn] (tr11x6) at (3,-1) {6};
\node[dn] (tr11x7) at (4,-1) {7};
\node[dn] (tr11x8) at (5,-1) {8};
\node[dn] (tr11x9) at (6,-1) {9};
\node[dn,dnhl] (tr11x10) at (7,-1) {10};
\draw[dptr] (tr11x1) -- (tr11x5);
\draw[dptr] (tr11x2) -- (tr11x1);
\draw[dptr] (tr11x3) -- (tr11x5);
\draw[dptr] (tr11x4) -- (tr11x5);
\draw[dptr] (tr11x6) -- (tr11x5);
\draw[dptr] (tr11x7) -- (tr11x5);
\draw[dptr] (tr11x8) -- (tr11x5);
\draw[dptr] (tr11x9) -- (tr11x5);
\draw[dptr] (tr11x10) -- (tr11x5);
\node[plab, anchor=north] at (3.5,-2.45) {(d) 第 11 步之後};
\end{tikzpicture}
\diagcaption{上表幾個時間點的森林。 藍框是根，橘色是這一步父指標有改的節點，箭頭由子
節點指向父節點。}
\label{fig:ds-trace}
\end{center}

第 1--4 步都是兩個單點合併，一樣大，$y$ 的根接到 $x$ 底下。 其他幾步要注意的地方：

\begin{itemize}
\item 第 5 步：\proc{Find-Set}$(2) = 1$，\proc{Find-Set}$(4) = 3$，兩棵都是 $2$ 個節點，
$3$ 接到 $1$ 底下，$\mathit{size}[1] = 4$。 注意改的是 $p[3]$，不是 $p[4]$。
\item 第 8 步：兩個根 $5$ 和 $1$ 都是 $4$ 個節點，$1$ 接到 $5$ 底下。 這時 $4$ 的深度
是 $3$，路徑 $4 \to 3 \to 1 \to 5$ 是整個森林最長的，圖 (b) 的樹高度是 $4$。
\item 第 9 步：\proc{Find-Set}$(4)$ 走過 $4, 3, 1, 5$，把 $4$ 和 $3$ 改接到 $5$；$1$ 本
來就是 $5$ 的子節點，不用改。 $2$ 不在這條路徑上，仍然掛在 $1$ 底下。
\item 第 10 步：\proc{Find-Set}$(10) = 9$，\proc{Find-Set}$(8)$ 走過 $8, 7, 5$，順便把
$8$ 改接到 $5$。 $\mathit{size}[9] = 2 < \mathit{size}[5] = 8$，兩個根交換，$9$ 接到
$5$ 底下。 這一步改了兩格：$p[8]$ 來自壓縮，$p[9]$ 來自合併。
\end{itemize}

\begin{examnote}
追蹤題通常會規定：(1) 一樣大（或秩相同）時誰當根；(2) \proc{Union}$(x, y)$ 裡的
\proc{Find-Set} 要不要壓縮；(3) 陣列的格式。 有些課本讓根的格子存\textbf{負的節點數}
（例如 $p[1] = -4$ 表示 $1$ 是根、這棵樹有 $4$ 個節點），不另外開 $\mathit{size}$
陣列；照這種寫法，上表第 6 步之後是 $-4, 1, 1, 3, -4, 5, 5, 7, -1, -1$。 這一系的課本
把依大小合併叫\textbf{加權規則}（weighting rule），把路徑壓縮叫\textbf{塌陷規則}
（collapsing rule）。 題目沒有規定時，在答案開頭寫清楚自己的約定，再開始畫。
\end{examnote}

\section{成本}
\label{sec:ds-cost}

\subsection{依大小或依秩合併：高度是 $O(\lg n)$}

\begin{theorem}
\label{thm:ds-size-height}
只用依大小合併（有沒有路徑壓縮都可以），任何時候，一棵有 $s$ 個節點的樹，高度不超
過 $\lfloor \lg s \rfloor + 1$。 所以 \proc{Find-Set} 最壞是 $O(\lg n)$。
\end{theorem}

\begin{proof}
先看沒有路徑壓縮的情形。 固定一個節點 $x$，它的深度只有在一種情況會增加：$x$ 所在
的樹的根 $r_y$ 被接到另一個根 $r_x$ 底下，這時深度加一。 依大小合併保證
$\mathit{size}[r_x] \ge \mathit{size}[r_y]$，所以合併後 $x$ 所在的樹的節點數
$\mathit{size}[r_x] + \mathit{size}[r_y] \ge 2\,\mathit{size}[r_y]$，至少是原來的兩
倍。 $x$ 一開始在只有 $1$ 個節點的樹裡；深度變成 $d$ 時，它所在的樹至少有 $2^d$ 個
節點。 現在這棵樹有 $s$ 個節點，所以 $2^d \le s$，$d \le \lfloor \lg s \rfloor$。
樹的高度是最大深度加一，不超過 $\lfloor \lg s \rfloor + 1$。

路徑壓縮只會把節點改接到更上面的祖先，任何節點的深度都不會因此變大，每棵樹的節點數
也不變，所以上界仍然成立。
\end{proof}

依秩合併的上界要先整理秩的性質。

\begin{lemma}
\label{lem:ds-rank}
使用依秩合併（有沒有路徑壓縮都可以），任何時候：
\begin{enumerate}
\item 若 $x$ 不是根，$\mathit{rank}[x] < \mathit{rank}[p[x]]$。
\item 秩為 $r$ 的根，它的樹至少有 $2^r$ 個節點。 所以每個秩都不超過
$\lfloor \lg n \rfloor$。
\item 以 $x$ 為根的子樹，從 $x$ 往下最長的路徑不超過 $\mathit{rank}[x]$ 條邊；樹的高
度不超過根的秩加一，也就不超過 $\lfloor \lg n \rfloor + 1$。
\end{enumerate}
\end{lemma}

\begin{whybox}[引理~\ref{lem:ds-rank} 的證明]
(1) $x$ 變成子節點的那一刻（虛擬碼~\ref{alg:ds-rank} 第 8 行），$\mathit{rank}[x] \le
\mathit{rank}[p[x]]$，相等時第 10 行把父節點的秩加一，所以嚴格小於。 之後 $x$ 不再是
根，秩不再變；它的父節點如果還是根，秩只會變大；路徑壓縮把 $p[x]$ 換成更上面的祖先，
由同樣的道理，祖先的秩更大。 不等式一直成立。

(2) 對操作的次數做歸納。 一開始每個根的秩是 $0$，樹有 $1 = 2^0$ 個節點。 合併時，
如果兩個根的秩不同，留下來的根秩不變、節點只會變多；如果兩個根的秩都是 $r$，各至少
$2^r$ 個節點，合併後的根秩是 $r + 1$，節點至少 $2^{r+1}$ 個。 路徑壓縮不改秩，也不
改每棵樹的節點數。 一棵樹最多 $n$ 個節點，$2^r \le n$，得到 $r \le \lfloor \lg n
\rfloor$。 不是根的節點，秩停在它最後一次當根時的值，也不超過這個上界。

(3) 從 $x$ 往下走，由 (1)，每走一條邊秩至少減一，而秩不小於 $0$，所以最多走
$\mathit{rank}[x]$ 條邊。 樹的高度（數節點）是最長路徑的邊數加一，不超過
$\mathit{rank}[\text{根}] + 1 \le \lfloor \lg n \rfloor + 1$。
\end{whybox}

這個上界做得到。 $n = 8$ 時依序做 \proc{Union}$(1, 2)$、$(3, 4)$、$(5, 6)$、
$(7, 8)$、$(1, 3)$、$(5, 7)$、$(1, 5)$，每次都是兩棵一樣的樹合併：根 $1$ 的秩是 $3$，高度 $4 = \lg 8 + 1$，$p[1..8] = 1, 1,
1, 3, 1, 5, 5, 7$。 依大小合併在這串操作上做出同一棵樹。 這時每次
\proc{Find-Set}$(8)$ 都要經過 $4$ 個節點（不壓縮的話），一般來說，只用其中一種合併規
則，$m$ 個操作的總時間是 $\Theta(m \lg n)$。


\subsection{加上路徑壓縮：$O(m\,\alpha(n))$}

兩個技巧一起用，總時間的上界降到幾乎是線性。 這裡只給結論，證明要用到精細的平攤分
析（位能法，第~3~章），超出本書的範圍。

\begin{theorem}[Tarjan 1975；CLRS 第 4 版 19.4 節]
\label{thm:ds-alpha}
森林表示同時使用依秩合併和路徑壓縮時，$m$ 個操作（其中 $n$ 個是 \proc{Make-Set}）
最壞情況的總時間是 $O(m\,\alpha(n))$。
\end{theorem}

$\alpha(n)$ 是一個成長極慢的函數，是 Ackermann 函數的某種反函數。 本書用 CLRS 的版
本：對 $k \ge 0$、$j \ge 1$，令
\[
A_0(j) = j + 1, \qquad A_k(j) = \underbrace{A_{k-1}(A_{k-1}(\cdots A_{k-1}}_{j + 1
\text{ 次}}(j)\cdots)) \quad (k \ge 1),
\]
再定義 $\alpha(n) = \min\{k : A_k(1) \ge n\}$。 算幾個值：$A_1(j) = 2j + 1$，
$A_2(j) = 2^{j+1}(j + 1) - 1$，所以 $A_0(1) = 2$、$A_1(1) = 3$、$A_2(1) = 7$、$A_3(1)
= A_2(A_2(1)) = A_2(7) = 2047$，而 $A_4(1) = A_3(2047)$ 已經遠大於 $10^{80}$。 下表把
$\alpha(n)$ 和 $\lg n$、迭代對數 $\lg^* n$（$n$ 要連續取幾次 $\lg$ 才會不超過 $1$）並
排：

\begin{center}
\begin{tabular}{rrrr}
\toprule
$n$ & $\lg n$ & $\lg^* n$ & $\alpha(n)$ \\
\midrule
$10$ & $3.3$ & $3$ & $3$ \\
$10^3$ & $10.0$ & $4$ & $3$ \\
$10^6$ & $19.9$ & $5$ & $4$ \\
$10^{9}$ & $29.9$ & $5$ & $4$ \\
$10^{18}$ & $59.8$ & $5$ & $4$ \\
$2^{65536}$ & $65536$ & $5$ & $4$ \\
\bottomrule
\end{tabular}
\end{center}

任何實際能存進電腦的 $n$，$\alpha(n) \le 4$。 所以定理~\ref{thm:ds-alpha} 實際上就是
「平攤每個操作常數時間」，但在數學上 $\alpha(n)$ 會無限變大，不能寫成 $O(m)$。
Tarjan 原本的結果用的是兩個參數的 $\alpha(m, n)$，CLRS 第 4 版改用單參數的
$\alpha(n)$；對任何實際的 $m$、$n$，兩者都不超過 $4$。 更早，Hopcroft 和 Ullman（1973）證明了 $O(m \lg^* n)$ 的上界，也很
接近線性。

幾個相關的結果，一樣只列結論：
\begin{itemize}
\item 依大小合併換掉依秩合併，或路徑壓縮換成路徑減半、路徑分割，上界一樣是
$O(m\,\alpha(n))$（Tarjan 與 van Leeuwen 1984）。
\item 只用路徑壓縮、不用任何合併規則時，$n$ 個 \proc{Make-Set}、最多 $n - 1$ 個
\proc{Union} 和 $f$ 個 \proc{Find-Set} 的最壞總時間是
$\Theta\bigl(n + f \cdot (1 + \log_{2 + f/n} n)\bigr)$（同一篇論文）。
\item 在相當一般的計算模型（cell probe 模型）下，任何互斥集合的資料結構都要 $\Omega(m\,\alpha(n))$ 的
時間（Fredman 與 Saks 1989），所以這個上界已經無法再改進。
\end{itemize}

最後要分清楚\textbf{平攤}和\textbf{單次最壞}。 $O(m\,\alpha(n))$ 是 $m$ 個操作的總
和，平均分給每個操作才是 $O(\alpha(n))$。 單獨一次 \proc{Find-Set} 仍然可能走過一
條長度 $\lfloor \lg n \rfloor + 1$ 的路徑（引理~\ref{lem:ds-rank}），例如上面
$8$ 個節點的例子第一次找 $8$ 的時候。 各種做法整理如下。

\begin{center}
\small
\begin{tabular}{p{0.30\linewidth}p{0.24\linewidth}p{0.36\linewidth}}
\toprule
做法 & 單次最壞 & $m$ 個操作的總時間 \\
\midrule
串列，不加權 & \proc{Find} $O(1)$，\proc{Union} $O(n)$ & $\Theta(n^2)$ 做得到（$m = 2n - 1$） \\
串列，加權合併 & \proc{Find} $O(1)$，\proc{Union} $O(n)$ & $O(m + n \lg n)$ \\
森林，不用技巧 & $O(n)$ & $O(mn)$ \\
依大小或依秩合併 & $O(\lg n)$ & $O(m \lg n)$ \\
只用路徑壓縮 & $O(n)$ & 見上面第二點 \\
依大小或依秩＋路徑壓縮（或減半、分割） & $O(\lg n)$ & $O(m\,\alpha(n))$ \\
\bottomrule
\end{tabular}
\end{center}

\section{C 實作}
\label{sec:ds-c}

C 程式用陣列存森林：\texttt{parent[x]} 是父節點，\texttt{size[x]} 只在根上有意
義。 元素從 $0$ 編號；題目或輸入從 $1$ 編號時，讀進來先減一，印出去再加一，整段一起
換，不在演算法裡面加減。

\begin{cbookrecap}[第~14~章 14.5 節、第~15~章 15.4 節、第~19~章 19.3 節]
\texttt{DSU *d} 是指向結構的指標，\texttt{d->parent} 等於 \texttt{(*d).parent}。 結
構裡的兩個陣列用 \texttt{malloc} 要記憶體，所以結構本身只存兩個指標；用完要各
\texttt{free} 一次。 第一個 \texttt{malloc} 成功、第二個失敗時，第一塊也要放掉，
\texttt{free(NULL)} 什麼都不做，所以可以直接兩個都 \texttt{free}。

遞迴函式 \texttt{find\_rec} 每呼叫一層就在呼叫堆疊上多一個框，路徑多長就疊幾層。
\end{cbookrecap}

下面的程式實作依大小合併，以及遞迴和迴圈兩種路徑壓縮：\texttt{dsu\_union} 裡用迴
圈版，單獨的 \texttt{F} 操作用遞迴版。 兩個版本都把整條路徑接到根，結果一樣。 程式
照第~\ref{sec:ds-trace}~節的 $11$ 個操作執行，每一步印出 \texttt{parent} 陣列。

\begin{example}[label=ex:ds-trace]{c}{(dsu\_trace.c) 依大小合併與路徑壓縮，印出每一步的父節點陣列}
#include <stdio.h>
#include <stdlib.h>

typedef struct {
    int n;          // 元素是 0..n-1
    int *parent;    // parent[x] == x 表示 x 是根
    int *size;      // 只有根的 size 有意義：這棵樹有幾個節點
    int sets;       // 目前有幾個集合
} DSU;

int dsu_init(DSU *d, int n)
{
    d->parent = malloc(n * sizeof *d->parent);
    d->size = malloc(n * sizeof *d->size);
    if (d->parent == NULL || d->size == NULL) {
        free(d->parent);
        free(d->size);
        return 0;
    }
    for (int x = 0; x < n; x++) {    // n 次 Make-Set
        d->parent[x] = x;
        d->size[x] = 1;
    }
    d->n = n;
    d->sets = n;
    return 1;
}

void dsu_free(DSU *d)
{
    free(d->parent);
    free(d->size);
}

// 路徑壓縮，遞迴版：回來的路上把每個節點直接接到根。
int find_rec(int *parent, int x)
{
    if (parent[x] != x) {
        parent[x] = find_rec(parent, parent[x]);
    }
    return parent[x];
}

// 路徑壓縮，迴圈版：第一趟找根，第二趟把路上的節點接到根。
int find_iter(int *parent, int x)
{
    int root = x;
    while (parent[root] != root) {
        root = parent[root];
    }
    while (parent[x] != root) {
        int next = parent[x];
        parent[x] = root;
        x = next;
    }
    return root;
}

// 依大小合併，一樣大時 y 的根接到 x 的根底下。
int dsu_union(DSU *d, int x, int y)
{
    int rx = find_iter(d->parent, x);
    int ry = find_iter(d->parent, y);
    if (rx == ry) {
        return 0;
    }
    if (d->size[rx] < d->size[ry]) {   // 讓 rx 是比較大的那棵
        int t = rx;
        rx = ry;
        ry = t;
    }
    d->parent[ry] = rx;
    d->size[rx] += d->size[ry];
    d->sets--;
    return 1;
}

// 印出 parent 陣列。 程式從 0 編號，印的時候加 1。
void print_parent(const DSU *d)
{
    for (int x = 0; x < d->n; x++) {
        printf(x > 0 ? " %d" : "%d", d->parent[x] + 1);
    }
    printf("\n");
}

int main(void)
{
    // 手算追蹤的 11 個操作，電腦從 1 編號
    struct { char op; int a, b; } ops[] = {
        {'U', 1, 2}, {'U', 3, 4}, {'U', 5, 6}, {'U', 7, 8},
        {'U', 2, 4}, {'U', 6, 8}, {'U', 9, 10}, {'U', 6, 2},
        {'F', 4, 0}, {'U', 10, 8}, {'F', 10, 0},
    };
    DSU d;
    if (!dsu_init(&d, 10)) {
        return 1;
    }
    for (int i = 0; i < 11; i++) {
        int a = ops[i].a - 1, b = ops[i].b - 1;
        if (ops[i].op == 'U') {
            dsu_union(&d, a, b);
            printf("U %d %d: ", a + 1, b + 1);
        } else {
            int r = find_rec(d.parent, a);
            printf("F %d -> %d: ", a + 1, r + 1);
        }
        print_parent(&d);
    }
    printf("sets = %d\n", d.sets);
    dsu_free(&d);
    return 0;
}
\end{example}
\iosample{}{U 1 2: 1 1 3 4 5 6 7 8 9 10 \\ U 3 4: 1 1 3 3 5 6 7 8 9 10 \\ U 5 6: 1 1 3 3 5 5 7 8 9 10 \\ U 7 8: 1 1 3 3 5 5 7 7 9 10 \\ U 2 4: 1 1 1 3 5 5 7 7 9 10 \\ U 6 8: 1 1 1 3 5 5 5 7 9 10 \\ U 9 10: 1 1 1 3 5 5 5 7 9 9 \\ U 6 2: 5 1 1 3 5 5 5 7 9 9 \\ F 4 -> 5: 5 1 5 5 5 5 5 7 9 9 \\ U 10 8: 5 1 5 5 5 5 5 5 5 9 \\ F 10 -> 5: 5 1 5 5 5 5 5 5 5 5 \\ sets = 1}

每一列都和第~\ref{sec:ds-trace}~節的表相同。 \texttt{dsu\_union} 先 find 再比大小，
所以合併時的兩次 find 也會壓縮路徑，第 10 步的 \texttt{p[8]} 就是這樣改的。 只問
「有幾組」時不必掃描整個陣列：每次真正合併讓 \texttt{sets} 減一就好。

\subsection{應用：拔線之後還剩幾個網段}
\label{sec:ds-app}

互斥集合只會合併，不會拆開。 如果實驗室要做斷線演練：先接好 $m$ 條線，再照一張清單
依序拔掉其中 $k$ 條，每拔一條就要回報還剩幾個網段、最大的網段有幾台電腦，該怎麼辦？

清單是事先就知道的，所以可以\textbf{倒過來做}。 拔完全部 $k$ 條之後剩下的線，是
「從來沒被拔過」的那些；先把它們接上，得到最後的狀態。 接著從清單的最後一條往前，
一條一條接回去：接回第 $i$ 條之前的狀態，正好是「依序拔完第 $1$ 到第 $i$ 條」之後
的狀態。 倒著做只有合併，互斥集合就用得上了。 最大的網段也一樣：倒著做時集合只會
變大，每次合併後和目前的最大值比一下就好。

\begin{example}[label=ex:ds-lab-cut]{c}{(lab\_cut.c) 依序拔線，倒過來用互斥集合回答}
#include <stdio.h>
#include <stdlib.h>

typedef struct {
    int *parent;
    int *size;
    int sets;       // 目前有幾個網段
    int largest;    // 目前最大的網段有幾台
} DSU;

typedef struct {
    int u, v;       // 兩端的電腦（程式裡從 0 編號）
    int cut;        // 1：這條線會被拔掉
} Cable;

typedef struct {
    int sets, largest;
} State;

int find(int *parent, int x)        // 迴圈版的路徑壓縮
{
    int root = x;
    while (parent[root] != root) {
        root = parent[root];
    }
    while (parent[x] != root) {
        int next = parent[x];
        parent[x] = root;
        x = next;
    }
    return root;
}

void unite(DSU *d, int x, int y)    // 依大小合併
{
    int rx = find(d->parent, x);
    int ry = find(d->parent, y);
    if (rx == ry) {
        return;
    }
    if (d->size[rx] < d->size[ry]) {
        int t = rx;
        rx = ry;
        ry = t;
    }
    d->parent[ry] = rx;
    d->size[rx] += d->size[ry];
    d->sets--;
    if (d->size[rx] > d->largest) {     // 最大值只會變大
        d->largest = d->size[rx];
    }
}

int main(void)
{
    int n, m, k = 0;
    if (scanf("%d %d", &n, &m) != 2 || n <= 0 || m < 0) {
        return 1;
    }
    Cable *c = malloc((m + 1) * sizeof *c);       // 線 1..m
    int *order = malloc((m + 1) * sizeof *order); // 拔的順序
    State *ans = malloc((m + 1) * sizeof *ans);
    DSU d;
    d.parent = malloc(n * sizeof *d.parent);
    d.size = malloc(n * sizeof *d.size);
    int ok = c != NULL && order != NULL && ans != NULL
             && d.parent != NULL && d.size != NULL;
    for (int e = 1; ok && e <= m; e++) {
        ok = scanf("%d %d", &c[e].u, &c[e].v) == 2;
        c[e].u--;
        c[e].v--;
        c[e].cut = 0;
    }
    ok = ok && scanf("%d", &k) == 1 && k >= 0 && k <= m;
    for (int i = 1; ok && i <= k; i++) {
        ok = scanf("%d", &order[i]) == 1
             && order[i] >= 1 && order[i] <= m;
        if (ok) {
            c[order[i]].cut = 1;
        }
    }
    if (ok) {
        for (int x = 0; x < n; x++) {
            d.parent[x] = x;
            d.size[x] = 1;
        }
        d.sets = n;
        d.largest = 1;
        for (int e = 1; e <= m; e++) {      // 先接從沒拔的線
            if (!c[e].cut) {
                unite(&d, c[e].u, c[e].v);
            }
        }
        for (int i = k; i >= 1; i--) {      // 倒著把線接回去
            ans[i].sets = d.sets;       // 拔完第 i 條之後
            ans[i].largest = d.largest;
            unite(&d, c[order[i]].u, c[order[i]].v);
        }
        printf("cut ends groups largest\n");
        printf("- - %d %d\n", d.sets, d.largest);   // 還沒拔線
        for (int i = 1; i <= k; i++) {
            Cable *e = &c[order[i]];
            printf("%d %d-%d %d %d\n", order[i], e->u + 1, e->v + 1,
                   ans[i].sets, ans[i].largest);
        }
    }
    free(c);
    free(order);
    free(ans);
    free(d.parent);
    free(d.size);
    return ok ? 0 : 1;
}
\end{example}
\iosample{8 9 \\ 1 2 \\ 2 3 \\ 3 1 \\ 3 4 \\ 5 6 \\ 6 7 \\ 4 5 \\ 7 8 \\ 6 8 \\ 5 \\ 2 7 9 6 1}{cut ends groups largest \\ - - 1 8 \\ 2 2-3 1 8 \\ 7 4-5 2 4 \\ 9 6-8 2 4 \\ 6 6-7 3 4 \\ 1 1-2 4 3}

輸入第一列是電腦數和線數，接著 $m$ 條線（依序是第 $1$ 到第 $m$ 條），再來是要拔幾
條、依序拔哪幾條；假設每條線最多拔一次。 輸出第二列是還沒拔線時的狀態。 拔第 $2$ 條
（2--3）時，$1, 2, 3$ 之間還有另外兩條線，網段數不變；拔第 $7$ 條（4--5）把教室切成
兩半；第 $9$ 條（6--8）拔掉後 $6, 7, 8$ 仍然繞得過去。 每個問題只多花一次合併，總時
間是 $O((m + n)\,\alpha(n))$。 如果拔線的順序要邊做邊決定（線上的版本），這招就不能
用了，需要更複雜的資料結構，不在本書的範圍。

\section{本章常見錯誤}
\label{sec:ds-pitfalls}

\begin{center}
\begin{tabular}{p{0.24\linewidth}p{0.50\linewidth}p{0.16\linewidth}}
\toprule
錯誤 & 例子 & 見 \\
\midrule
接的不是根 & \proc{Union} 寫成 $p[y] \gets x$；$y$ 原來的根和其他節點留在舊的樹裡 &
第~\ref{sec:ds-forest}~節 \\
比父節點不比代表 & 判斷同一組時比 $p[a] = p[b]$，兩個節點在同一棵樹的不同深度就答錯 &
第~\ref{sec:ds-list}~節 \\
大小重複累加 & $x$、$y$ 已經同一組還照樣 $\mathit{size}[r_x] \gets \mathit{size}[r_x]
+ \mathit{size}[r_y]$；或把大小加在非根節點上 & 第~\ref{sec:ds-size}~節 \\
把秩當成高度 & 路徑壓縮後去更新秩；用實際樹高決定誰當根；忘了秩相等時加一，或秩不等也
加一 & 第~\ref{sec:ds-rank}~節 \\
高度差一 & 說「秩 $r$ 的樹高度是 $r$」，沒有換成本書數節點的定義（高度 $\le r + 1$） &
第~\ref{sec:ds-rank}~節 \\
\proc{Union} 裡的 find 不壓縮 & 手算時只在單獨的 \proc{Find-Set} 才壓縮，漏改了合併時
經過的節點 & 第~\ref{sec:ds-trace}~節 \\
平攤當成單次 & 說「每次 \proc{Find-Set} 最壞是 $O(\alpha(n))$」 & 第~\ref{sec:ds-cost}~節 \\
遞迴太深 & 沒有合併規則時用遞迴的 find，鏈長 $10^6$ 就用光堆疊 & 第~\ref{sec:ds-compress}~節 \\
以為可以拆開 & 想用互斥集合直接處理刪除線路；應該倒過來做，或改用別的方法 &
第~\ref{sec:ds-app}~節 \\
\bottomrule
\end{tabular}
\end{center}

\section{觀念題}
\label{sec:ds-exercises}

題型和研究所考試、期中期末考的筆試相近，有選擇、計算、追蹤和證明。 沒有特別說明時，
\proc{Union}$(x, y)$ 在兩個根一樣大（或秩相同）時讓 $y$ 的根接到 $x$ 的根底下，
\proc{Union} 裡的 \proc{Find-Set} 也照題目規定的方式壓縮；高度數節點。 解答在本章最
後。

\begin{enumerate}
\item （選擇，可能不只一個答案）下列哪些敘述正確？
\begin{enumerate}
\item 依秩合併、不做路徑壓縮時，每個節點的秩等於從它往下最長路徑的邊數。
\item 做路徑壓縮時，要把路徑上節點的秩改成新的值，否則之後的合併會選錯根。
\item 只用依大小合併時，$100$ 個元素的森林裡，最高的樹高度不超過 $7$。
\item 依秩合併加路徑壓縮時，每一次 \proc{Find-Set} 最壞都是 $O(\alpha(n))$。
\item 串列表示用加權合併時，\proc{Find-Set} 是 $O(1)$。
\end{enumerate}

\item （計算）用串列表示，\proc{Make-Set} 了 $1$ 到 $8$ 之後依序做
\proc{Union}$(1, 2)$、$(3, 1)$、$(4, 1)$、$(5, 6)$、$(5, 7)$、$(8, 5)$、$(8, 4)$。
分別在 (a) 一律把 $y$ 的串列接到 $x$ 的後面、(b) 加權合併（一樣長時接 $y$ 的）兩種
做法下，求每次合併改了幾個 $\mathit{set}$ 指標和總數，以及最後的代表。

\item \label{q:ds-rank-trace}（追蹤）$9$ 個元素，用依秩合併（虛擬碼~\ref{alg:ds-rank}）
加路徑壓縮，依序做 \proc{Union}$(1, 2)$、$(3, 4)$、$(4, 2)$、$(5, 6)$、$(7, 5)$、
$(8, 9)$、$(9, 6)$、$(4, 7)$，最後 \proc{Find-Set}$(6)$。 寫出每一步之後的 $p[1..9]$
和 $\mathit{rank}[1..9]$。 最後一步之後，根的秩和樹的高度各是多少？ 這和
第~\ref{sec:ds-rank}~節說的關係一致嗎？

\item （追蹤）用「根存負的節點數」的陣列（第~\ref{sec:ds-trace}~節作答提醒），目前
\[
p[1..8] = -3,\ 1,\ 1,\ -4,\ 4,\ 5,\ 4,\ -1.
\]
依序做 \proc{Union}$(3, 6)$、\proc{Find-Set}$(2)$、\proc{Union}$(8, 2)$，用加權規則和塌
陷規則（依大小合併、路徑壓縮）。 寫出每一步之後的陣列。

\item \label{q:ds-compare}（比較）$9$ 個元素，依序做 \proc{Union}$(2, 1)$、$(3, 1)$、$(4, 3)$、
$(2, 9)$、$(6, 5)$、$(7, 8)$、$(6, 8)$、$(5, 4)$，最後 \proc{Find-Set}$(8)$。
分別用下列五種做法，寫出最後的 $p[1..9]$、最高的樹的高度，以及最後的
\proc{Find-Set}$(8)$ 經過幾個節點。
(a) 不用任何技巧（虛擬碼~\ref{alg:ds-naive}）；(b) 依大小合併，不壓縮；(c) 依秩合
併，不壓縮；(d) 依大小合併＋路徑壓縮；(e) 依秩合併＋路徑壓縮。
再說明 (b) 和 (c) 在第 8 個操作為什麼選了不同的根。

\item （證明）只用依大小合併、不壓縮。 對合併的次數做歸納，證明：高度是 $h$ 的樹至
少有 $2^{h-1}$ 個節點。 這和定理~\ref{thm:ds-size-height} 的證明有什麼不同？

\item （證明）依秩合併（有沒有路徑壓縮都可以），證明：在任何時候，秩恰好是 $r$ 的節
點最多 $n / 2^r$ 個。（提示：把每個秩 $r$ 的節點，對應到它變成秩 $r$ 那一刻樹裡的
節點。）

\item （追蹤）森林的父節點是 $p[2] = 1$、$p[3] = 1$、$p[4] = 3$、$p[5] = 3$、
$p[6] = 5$、$p[7] = 5$、$p[8] = 7$、$p[9] = 7$，$1$ 是根。 分別用路徑壓縮、路徑減
半、路徑分割做一次 \proc{Find-Set}$(9)$，寫出之後的 $p[1..9]$ 和樹的高度。

\item （計算）
\begin{enumerate}
\item 求 $\alpha(5)$、$\alpha(100)$、$\alpha(5000)$、$\alpha(10^{30})$。
\item 求 $\lg^* 16$、$\lg^* 65536$、$\lg^* 2^{65536}$。
\item 「依秩合併＋路徑壓縮的總時間是 $O(m\,\alpha(n))$，而 $\alpha(n) \le 4$，所以總
時間是 $O(m)$。」 這句話哪裡不對？
\end{enumerate}

\item （追蹤）範例~\ref{ex:ds-lab-cut} 的線路不變，把拔線的清單改成 \texttt{4} 條：
\texttt{7 3 2 9}。 (a) 程式倒著接回去的順序是什麼？ (b) 寫出程式的輸出。 (c) 為什
麼互斥集合不能直接照拔線的順序處理？
\end{enumerate}

\section{實作練習}
\label{sec:ds-practice}

下面是 Aki-Judge 上和本章相關的題目。 題目裡的「不相交集合」「併查集」就是本章的互斥
集合。 交之前先用第~\ref{sec:ds-trace}~節的操作序列和範例~\ref{ex:ds-trace} 的輸出自己
對一次。

\begin{center}
\begin{tabular}{llc}
\toprule
小節 & 題目 & 難度 \\
\midrule
應用 & \akiref{counting-connected-components}{計算連通分量數} & 暖身 \\
森林表示 & \akiref{disjoint-set-union-data-structure}{不相交集合資料結構} & 基本 \\
 & \akiref{merges-and-connectivity-queries}{合併與連通查詢} & 基本 \\
依大小合併 & \akiref{size-of-the-largest-component}{最大連通分量的大小} & 基本 \\
 & \akiref{device-group-size}{裝置群組大小} & 進階 \\
綜合 & \akiref{power-grid-regions}{配電網路分區} & 進階 \\
\bottomrule
\end{tabular}
\end{center}

幾題的提示：

\begin{itemize}
\item 「計算連通分量數」：每次真正的合併讓連通元件數減一，或是全部做完以後數有幾個
$i$ 滿足 $\proc{Find-Set}(i) = i$。 不要用 $n - m$，重複的邊和自環不會減少分量。
\item 「不相交集合資料結構」只交標頭檔和函式。 \texttt{init} 可能出現不只一次，每次
都要把整個結構重設；\texttt{count} 就是範例~\ref{ex:ds-trace} 的 \texttt{sets}。
\item 「合併與連通查詢」的操作多達 $5 \times 10^5$ 個，兩個技巧都要用。 比較的是兩個
\proc{Find-Set} 的結果，不是兩個父節點。
\item 「最大連通分量的大小」照第~\ref{sec:ds-size}~節，大小只記在根上，已經同一組時不
能再加。
\item 「裝置群組大小」的編號可以大到 $10^9$，不能直接當陣列索引。 先把所有出現過的
編號排序、去掉重複，再用二分搜尋把編號換成 $0, 1, 2, \ldots$（排序見第~14~章，二分搜尋可以用
第~1~章觀念題的 \proc{Lower-Bound}），之後就是一般的依大小合併。 $a = b$ 時也要回答所在
群組的大小。
\item 「配電網路分區」可以拆電纜，互斥集合不能拆。 站數最多 $30$，每次查詢時把目前所
有的電纜重新做一次合併就夠快；這是「資料小、但有刪除」時的直接做法。 電纜總長要另外
加總，每條只算一次。
\end{itemize}

\section{讀完你應該會}
\label{sec:ds-checklist}

\begin{itemize}
\item 說出互斥集合的三個操作和參數 $n$、$m$ 的意思，用它寫出「依序接線、隨時查詢」
這類程序。
\item 用串列表示做合併，數出改了幾個代表指標，並證明加權合併的總時間是 $O(m + n \lg
n)$。
\item 給一串操作和規定（依大小或依秩、平手規則、要不要壓縮），寫出每一步之後的父節點
陣列（含根存負數的寫法）和秩陣列，畫出森林。
\item 分清楚秩、高度、深度：單一節點的秩是 $0$、高度是 $1$；高度 $\le$ 秩 $+ 1$；路徑壓
縮不改秩。
\item 證明依大小或依秩合併的樹高度是 $O(\lg n)$，說出加上路徑壓縮後 $O(m\,\alpha(n))$
的結論與出處，並分清楚平攤和單次最壞。
\item 寫出 C 的陣列版互斥集合（遞迴和迴圈的路徑壓縮），知道遇到刪除時可以把操作倒過
來處理。
\end{itemize}

\textbf{延伸閱讀}（書末會整理完整的參考資料）

\begin{itemize}
\item Thomas H. Cormen 等，\textit{Introduction to Algorithms}，第 4 版，第 19 章：
串列與森林兩種表示、依秩合併與路徑壓縮，19.4 節用位能法完整證明 $O(m\,\alpha(n))$；
章末的附註整理了這個問題的研究歷史。
\item Robert E. Tarjan，``Efficiency of a Good But Not Linear Set Union Algorithm''，
\textit{Journal of the ACM} 22(2)，1975：第一次證明 $\alpha$ 的上界。
\item Robert E. Tarjan、Jan van Leeuwen，``Worst-Case Analysis of Set Union
Algorithms''，\textit{Journal of the ACM} 31(2)，1984：路徑減半、路徑分割等變形，以
及只用路徑壓縮時的上界。
\item Ellis Horowitz、Sartaj Sahni、Susan Anderson-Freed，\textit{Fundamentals of Data
Structures in C}，第 5 章：根存負數的陣列寫法、加權規則和塌陷規則，附 C 程式。
\end{itemize}

\addcontentsline{toc}{section}{觀念題解答}
\section*{觀念題解答}

\begin{enumerate}
\item (a)、(c)、(e) 正確。 (a)：不壓縮時，合併兩個秩 $r$ 的根正好讓最長路徑多一條邊，
秩不同時最長路徑不變，歸納可得。 (b)：路徑壓縮不改秩，秩本來就只是上界，不需要準確。
(c)：$\lfloor \lg 100 \rfloor + 1 = 6 + 1 = 7$。 (d)：$O(\alpha(n))$ 是平攤，單次最壞
可以是 $\Theta(\lg n)$。 (e)：讀 $x.\mathit{set}.\mathit{head}$ 就好。

\item 每次合併改的指標數：
\begin{center}
\begin{tabular}{lccccccc|c}
\toprule
合併 & $(1,2)$ & $(3,1)$ & $(4,1)$ & $(5,6)$ & $(5,7)$ & $(8,5)$ & $(8,4)$ & 合計 \\
\midrule
(a) 一律接 $y$ & 1 & 2 & 3 & 1 & 1 & 3 & 4 & 15 \\
(b) 加權 & 1 & 1 & 1 & 1 & 1 & 1 & 4 & 10 \\
\bottomrule
\end{tabular}
\end{center}
(a) 的最後一條串列是 $8, 5, 6, 7, 4, 3, 1, 2$，代表 $8$。 (b) 的 $(3, 1)$、$(4, 1)$、
$(8, 5)$ 都是把單一元素接到長串列後面；最後一次兩條都是 $4$ 個，接 $y$ 的，最後的串列
是 $5, 6, 7, 8, 1, 2, 3, 4$，代表 $5$。

\item 每列是做完那一步之後的狀態：
\begin{center}
\small
\begin{tabular}{ll}
\toprule
操作 & $p[1..9]$ \qquad $\mathit{rank}[1..9]$ \\
\midrule
\proc{Union}$(1, 2)$ & $1,1,3,4,5,6,7,8,9$ \qquad $1,0,0,0,0,0,0,0,0$ \\
\proc{Union}$(3, 4)$ & $1,1,3,3,5,6,7,8,9$ \qquad $1,0,1,0,0,0,0,0,0$ \\
\proc{Union}$(4, 2)$ & $3,1,3,3,5,6,7,8,9$ \qquad $1,0,2,0,0,0,0,0,0$ \\
\proc{Union}$(5, 6)$ & $3,1,3,3,5,5,7,8,9$ \qquad $1,0,2,0,1,0,0,0,0$ \\
\proc{Union}$(7, 5)$ & $3,1,3,3,5,5,5,8,9$ \qquad $1,0,2,0,1,0,0,0,0$ \\
\proc{Union}$(8, 9)$ & $3,1,3,3,5,5,5,8,8$ \qquad $1,0,2,0,1,0,0,1,0$ \\
\proc{Union}$(9, 6)$ & $3,1,3,3,8,5,5,8,8$ \qquad $1,0,2,0,1,0,0,2,0$ \\
\proc{Union}$(4, 7)$ & $3,1,3,3,8,5,8,3,8$ \qquad $1,0,3,0,1,0,0,2,0$ \\
\proc{Find-Set}$(6)$ & $3,1,3,3,3,3,8,3,8$ \qquad $1,0,3,0,1,0,0,2,0$ \\
\bottomrule
\end{tabular}
\end{center}
幾個關鍵：$(4, 2)$ 的兩個根 $3$、$1$ 秩都是 $1$，$1$ 接到 $3$ 底下，$3$ 的秩變 $2$。
$(7, 5)$ 是秩 $0$ 的 $7$ 接到秩 $1$ 的 $5$ 底下，秩都不變。 $(9, 6)$ 的兩個根 $8$、$5$
秩都是 $1$，$5$ 接到 $8$ 底下。 $(4, 7)$ 先 \proc{Find-Set}$(7)$，路徑 $7 \to 5 \to 8$
把 $7$ 改接到 $8$；兩個根 $3$、$8$ 秩都是 $2$，$8$ 接到 $3$ 底下，$3$ 的秩變 $3$。 這時
$6 \to 5 \to 8 \to 3$ 是最長的路徑，高度 $4 = 3 + 1$。 最後 \proc{Find-Set}$(6)$ 把 $6$、
$5$ 改接到 $3$，最長的路徑剩 $7 \to 8 \to 3$（或 $9 \to 8 \to 3$、$2 \to 1 \to 3$），
高度變成 $3$，根的秩仍然是 $3$。 高度 $3 \le 3 + 1$，符合「高度 $\le$ 秩 $+ 1$」；壓縮
以後等號不再成立。

\item 每一步之後的陣列：
\begin{center}
\begin{tabular}{ll}
\toprule
操作 & $p[1..8]$ \\
\midrule
開始 & $-3,\ 1,\ 1,\ -4,\ 4,\ 5,\ 4,\ -1$ \\
\proc{Union}$(3, 6)$ & $4,\ 1,\ 1,\ -7,\ 4,\ 4,\ 4,\ -1$ \\
\proc{Find-Set}$(2)$ & $4,\ 4,\ 1,\ -7,\ 4,\ 4,\ 4,\ -1$ \\
\proc{Union}$(8, 2)$ & $4,\ 4,\ 1,\ -8,\ 4,\ 4,\ 4,\ 4$ \\
\bottomrule
\end{tabular}
\end{center}
\proc{Union}$(3, 6)$：\proc{Find-Set}$(3) = 1$（$3$ 個節點），\proc{Find-Set}$(6)$ 走過
$6, 5, 4$，把 $6$ 改接到 $4$；$1$ 的樹比較小，接到 $4$ 底下，$p[4] = -(3 + 4) = -7$。
\proc{Find-Set}$(2)$ 走過 $2, 1, 4$，$2$ 改接到 $4$。 最後 $8$ 只有一個節點，接到 $4$
底下。 注意 $p[3] = 1$ 沒變：$3$ 從來不在被壓縮的路徑上。

\item 五種做法的結果：
\begin{center}
\small
\begin{tabular}{lccc}
\toprule
做法 & 最後的 $p[1..9]$ & 最高的樹高度 & \proc{Find-Set}$(8)$ 經過 \\
\midrule
(a) 不用技巧 & $2,3,4,6,6,6,6,7,4$ & $5$ & $3$（$8, 7, 6$）\\
(b) 依大小 & $2,2,2,2,6,2,6,7,2$ & $4$ & $4$（$8, 7, 6, 2$）\\
(c) 依秩 & $2,6,2,2,6,6,6,7,2$ & $3$ & $3$（$8, 7, 6$）\\
(d) 依大小＋壓縮 & $2,2,2,2,6,2,2,2,2$ & $3$ & $4$ \\
(e) 依秩＋壓縮 & $2,6,2,2,6,6,6,6,2$ & $3$ & $3$ \\
\bottomrule
\end{tabular}
\end{center}
(a) 每次都把 $y$ 的根接到 $x$ 的根底下：前三次做出鏈 $1 \to 2 \to 3 \to 4$，
\proc{Union}$(2, 9)$ 讓 $9$ 接到 $4$，最後 \proc{Union}$(5, 4)$ 讓 $4$ 接到 $6$，最長的
路徑是 $1 \to 2 \to 3 \to 4 \to 6$。 (b)、(c) 前七個操作的結果相同：$2$ 是 $5$ 個節點、
高度 $2$ 的星形（秩 $1$），$6$ 是 $4$ 個節點、高度 $3$ 的樹（秩 $2$）。 第 8 個操作依大小
選節點多的 $2$ 當根，高的樹被接到下面，高度變 $4$；依秩選秩大的 $6$ 當根，高度維持 $3$。
(d)、(e) 在前八個操作都沒有可以壓縮的路徑（每次 find 的節點都在根或根的子節點），所以
和 (b)、(c) 一樣，最後的 \proc{Find-Set}$(8)$ 才壓縮：(d) 把 $8$ 和 $7$ 接到 $2$，(e) 把
$8$ 接到 $6$。 兩個合併規則都保證 $O(\lg n)$ 的高度，但單看一串操作，誰比較矮不一定。

\item 歸納的對象是「目前森林裡的每棵樹都滿足：高度 $h$ 的樹至少 $2^{h-1}$ 個節點」。
一開始每棵樹只有一個節點，$h = 1$，$2^0 = 1$，成立。 合併兩棵樹 $T_x$（$s_x$ 個節點、
高度 $h_x$）和 $T_y$（$s_y$ 個節點、高度 $h_y$），設 $s_x \ge s_y$，$T_y$ 的根接到
$T_x$ 的根底下，新的高度是 $h = \max(h_x, h_y + 1)$。 若 $h = h_x$，節點數 $s_x + s_y
> s_x \ge 2^{h_x - 1}$。 若 $h = h_y + 1$，節點數 $s_x + s_y \ge 2 s_y \ge 2 \cdot
2^{h_y - 1} = 2^{h - 1}$。 其他樹不變，所以性質一直成立；$s \ge 2^{h-1}$ 等於 $h \le \lg
s + 1$。 定理~\ref{thm:ds-size-height} 是追蹤\textbf{單一節點}的深度什麼時候增加，這裡
是對\textbf{整棵樹}做歸納；兩種都常見，結論相同。

\item 對每個秩是 $r$ 的節點 $x$，看它的秩變成 $r$ 的那一刻：那時 $x$ 是根，由
引理~\ref{lem:ds-rank} (2)，它的樹至少有 $2^r$ 個節點，把這些節點記成 $N(x)$。 只要證明
不同的 $x$ 得到的 $N(x)$ 兩兩不相交，就有（秩 $r$ 的節點數）$\cdot\, 2^r \le n$。

假設 $z$ 同時在 $N(x)$ 和 $N(y)$ 裡，$x \ne y$，而且 $x$ 比 $y$ 先變成秩 $r$。 集合只會
合併、不會拆開，所以 $x$ 變成秩 $r$ 之後，$z$ 和 $x$ 一直在同一棵樹裡。 到了 $y$ 變成秩
$r$ 的那一刻，$z$ 在 $y$ 的樹裡，所以 $x$ 也在 $y$ 的樹裡，而 $y$ 是根、$x \ne y$，$x$
是 $y$ 的子孫。 由引理~\ref{lem:ds-rank} (1)，沿著 $x$ 到 $y$ 的路徑秩嚴格變大，
$\mathit{rank}[x] < \mathit{rank}[y] = r$；但秩不會變小，$\mathit{rank}[x] \ge r$，矛盾。
所以秩 $r$ 的節點最多 $n / 2^r$ 個。 （這也再一次說明秩不超過 $\lg n$。）

\item 路徑是 $9 \to 7 \to 5 \to 3 \to 1$。
\begin{center}
\begin{tabular}{lcc}
\toprule
做法 & 之後的 $p[1..9]$ & 高度 \\
\midrule
路徑壓縮 & $1,1,1,3,1,5,1,7,1$ & $3$ \\
路徑減半 & $1,1,1,3,1,5,5,7,5$ & $4$ \\
路徑分割 & $1,1,1,3,1,5,3,7,5$ & $4$ \\
\bottomrule
\end{tabular}
\end{center}
壓縮：$9, 7, 5$ 都改接到 $1$（$3$ 本來就是）；最長的路徑變成 $4 \to 3 \to 1$ 這類，高度
$3$。 減半：$9$ 改指祖父 $5$，跳到 $5$；$5$ 改指祖父 $1$，跳到 $1$ 結束；$7$ 沒有改。
分割：$9 \to 5$、$7 \to 3$、$5 \to 1$；$3$ 的父節點就是根，改指祖父還是根，不變。
減半之後還有 $8 \to 7 \to 5 \to 1$，分割之後還有 $8 \to 7 \to 3 \to 1$，都是 $4$ 個節點的
路徑，高度 $4$。

\item (a) $A_2(1) = 7 \ge 5$ 且 $A_1(1) = 3 < 5$，所以 $\alpha(5) = 2$。 $A_2(1) = 7 < 100 \le
A_3(1) = 2047$，$\alpha(100) = 3$。 $5000 > 2047$，$\alpha(5000) = 4$。 $10^{30} < A_4(1)$，
$\alpha(10^{30}) = 4$。

(b) $16 \to 4 \to 2 \to 1$，$\lg^* 16 = 3$；$65536 \to 16 \to 4 \to 2 \to 1$，
$\lg^* 65536 = 4$；$\lg^* 2^{65536} = 5$。

(c) $\alpha(n) \le 4$ 只對「實際會出現的 $n$」成立。 $O$ 描述的是 $n \to \infty$ 的行為，
$\alpha(n)$ 會無限變大，所以數學上不能寫成 $O(m)$；而且 Fredman 和 Saks 證明了
$\Omega(m\,\alpha(n))$ 的下界，總時間確實不是 $O(m)$。 在工程上說「幾乎線性」是可以的。

\item (a) 倒著接回去的順序是第 $9$、$2$、$3$、$7$ 條。 程式先接從沒被拔的第 $1, 4, 5,
6, 8$ 條：$\{1, 2\}$、$\{3, 4\}$、$\{5, 6, 7, 8\}$。
(b) 輸出：
\begin{center}
\ttfamily
\begin{tabular}{l}
cut ends groups largest \\
- - 1 8 \\
7 4-5 2 4 \\
3 3-1 2 4 \\
2 2-3 3 4 \\
9 6-8 3 4 \\
\end{tabular}
\end{center}
拔掉 4--5 把教室分成 $\{1, 2, 3, 4\}$ 和 $\{5, 6, 7, 8\}$；再拔 3--1，$1, 2, 3$ 之間還有
1--2、2--3，不變；拔 2--3 之後 $\{1, 2\}$ 和 $\{3, 4\}$ 分開；拔 6--8 不影響 $5$ 到 $8$。
(c) 拔線會讓一個集合分成兩個，互斥集合沒有拆開的操作；父節點陣列也沒有記錄是哪幾條線
把樹連起來的，拔掉一條線時無從知道該拆哪裡。
\end{enumerate}
